\documentclass[11pt,a4paper]{article}

% =========================================================================
% MLSM2154 - Gesture Recognition Project (Group 6)
% Internal pedagogical memo: WHAT the pipeline does, HOW it does it,
% and WHY each design decision was made, with literature backing.
%
% This document is NOT the final 10-page report for the professor. It is
% a complete reference written for ourselves, to memorise the pipeline
% before the oral defence and to never forget the rationale behind a
% given line of CODE_IA.py.
% =========================================================================

\usepackage[utf8]{inputenc}
\usepackage[T1]{fontenc}
\usepackage{lmodern}
\usepackage[a4paper,margin=2.5cm]{geometry}
\usepackage{amsmath,amssymb,amsthm}
\usepackage{graphicx}
\usepackage{booktabs}
\usepackage{xcolor}
\usepackage[colorlinks=true,linkcolor=blue!50!black,citecolor=blue!50!black,urlcolor=blue!50!black]{hyperref}
\usepackage{listings}
\usepackage{framed}
\usepackage{enumitem}
\usepackage{titlesec}
\usepackage{multicol}
\usepackage{caption}

\definecolor{whyboxbg}{RGB}{255,247,224}
\definecolor{whyboxbd}{RGB}{204,153,0}
\definecolor{codebg}{RGB}{246,246,246}

\lstset{
    basicstyle=\ttfamily\small,
    backgroundcolor=\color{codebg},
    frame=single,
    framerule=0pt,
    breaklines=true,
    columns=fullflexible,
    xleftmargin=0.5em,
    showstringspaces=false,
}

% Decision/Justification box
\newenvironment{whybox}[1]{%
    \begin{shaded*}\noindent\textbf{Pourquoi : #1}\par\vspace{0.3em}}
    {\end{shaded*}}

\renewenvironment{shaded*}
    {\definecolor{shadecolor}{named}{whyboxbg}\begin{shaded}}
    {\end{shaded}}

\titleformat{\section}
    {\Large\bfseries\color{blue!40!black}}{\thesection.}{0.6em}{}
\titleformat{\subsection}
    {\large\bfseries\color{blue!30!black}}{\thesubsection.}{0.6em}{}

\title{\textbf{Gesture Recognition Pipeline --- Documented Walk-Through}\\
\large MLSM2154 Group 6 internal pedagogical memo}
\author{Andry Lenny \and El Mohcine Mohamed \and Ottevaere Arthur}
\date{May 2026}

\begin{document}
\maketitle
\tableofcontents
\bigskip

% =========================================================================
\section{Goal of this document}

This file is not the final 10-page report we hand in to Prof.\ Saerens.
It is the long form of \emph{everything our code does and why}, written so
that any of us can re-read it the night before the oral defence and
recover the full mental model behind \texttt{CODE\_IA.py}.

The pipeline implements supervised classification of 3D hand-gesture
trajectories using two datasets (Domain~1 = numbers 0--9, Domain~4 = 3D
volumetric shapes, both from \cite{huang2019gesture}). For each gesture
we are given a sequence
$\bigl(\mathbf{r}(t)\bigr)_{t=1}^{T}\in\mathbb{R}^{3\times T}$ of 3D
positions.

Seven classifiers are compared:
\begin{enumerate}[itemsep=0.2em,topsep=0.2em]
    \item Edit Distance on $k$-means symbol sequences (baseline, DP).
    \item Dynamic Time Warping on the 3D points directly (baseline, DP).
    \item Decision Tree on hand-crafted features (advanced, pedagogical
          control).
    \item Random Forest on hand-crafted features (advanced).
    \item SVM with RBF kernel on hand-crafted features (advanced).
    \item Logistic Regression (multinomial, L2) on hand-crafted features
          (advanced, linear baseline).
    \item \$1 Recognizer adapted to 3D (advanced, template-based; evaluated
          on \emph{raw} data only --- see Sec.~\ref{sec:dollar-raw}).
\end{enumerate}

% =========================================================================
\section{Pipeline overview}

\begin{lstlisting}[caption={Phases as they appear in CODE\_IA.py}]
load -> standardise (per-gesture) -> PCA denoise 3D->2D->3D
    -> k-means VQ (per-fold)
    -> validation curves for K_CLUSTERS and KNN_K  (KNN_K forced to 1)
    -> ablation: 7 methods x 3 preprocessings, run in UI AND UD
       ($1 forced to raw afterwards per Wobbrock 2007/Kratz & Rohs 2010)
    -> per-fold permutation-importance feature selection
       (DT, RF, SVM, LR; cumulative threshold 90%)
    -> statistical tests: Wilcoxon n=100, Bonferroni + BH-FDR (21 pairs)
    -> confusion matrix of the best model per domain
    -> overfitting diagnostic: train-vs-test gap + learning curves
    -> summary CSV under Outputs/tables/summary/
\end{lstlisting}

The single Python file \texttt{CODE\_IA.py} is organised in 13 numbered
sections. Outputs are written under \texttt{Outputs/} in nested
sub-folders (see Sec.~\ref{sec:outputs}).

% =========================================================================
\section{Data loading}

\subsection{Domain 1}
CSV files \texttt{Subject\textit{i}-\textit{g}-\textit{r}.csv} with
columns \verb|x,y,z,t|. Filename regex
\verb|Subject(\d+)-(\d+)-(\d+)\.csv| gives subject, gesture, repetition.

\subsection{Domain 4}
Plain-text files (no extension) with a small header
\verb|Class id = ..., User id = ...| followed by a \verb|<x>| marker
and the 3D coordinates.

\subsection{Subject re-indexing to 0}
Both loaders apply \verb|subject - 1| so that user IDs are
zero-indexed (Pythonic). \textbf{This must match the assumptions made
in the cross-validation indexing.}

\subsection{Completeness check}
\verb|check_completeness| asserts that every (user, gesture) cell
contains exactly 10 repetitions. A warning is printed if not.

% =========================================================================
\section{Per-gesture standardisation}

For each gesture and each axis $a\in\{x,y,z\}$:
\[
    \tilde{r}_a(t) \;=\; \frac{r_a(t) - \mu_a^{\,(\text{gesture})}}
                              {\sigma_a^{\,(\text{gesture})}+\varepsilon}.
\]

\begin{whybox}{per-gesture standardisation}
A global mean/std across the whole dataset would normalise away the
\emph{class}-level differences between gestures. Per-gesture
standardisation removes scale/offset due to where in space the user
started drawing and how big the strokes were, while preserving the
\emph{shape} of the trajectory. This is the standard normalisation
for time-series classification \cite{theodoridis2009pattern}, and is
also the form expected by the \$1 Recognizer family before its rotation
and scaling steps.
\end{whybox}

% =========================================================================
\section{Per-gesture PCA denoising 3D $\to$ 2D $\to$ 3D}

For each gesture (a $T\times 3$ matrix $X$):
\begin{enumerate}[itemsep=0.2em,topsep=0.2em]
    \item Fit PCA with 3 components on the $T$ time-step points of this
          gesture only.
    \item Project onto the top 2 principal components (the trajectory's
          dominant plane in 3D).
    \item Set the third coordinate to zero (the orthogonal-to-plane
          component is treated as noise).
    \item Back-project to 3D using the inverse PCA transform.
    \item Save the 3 explained-variance ratios (EVR) as additional
          features for the tree- and kernel-based classifiers.
\end{enumerate}

\begin{whybox}{denoise in the plane of the gesture}
Most digits in Domain~1 are essentially planar (sketched in mid-air
on a roughly vertical plane), so most of the variance is captured by
PC1 + PC2. Removing PC3 removes hand jitter that is orthogonal to
that plane \cite{jolliffe2002pca}. We keep the EVR triple because it
encodes \emph{how planar} the gesture was, which is a discriminative
feature for Domain~4 (cones and tori are intrinsically 3D and have
non-trivial PC3).
\end{whybox}

% =========================================================================
\section{$k$-means vector quantisation (for Edit Distance)}

Edit Distance requires a discrete sequence of symbols, not 3D points.
For each fold of the user-independent CV:
\begin{enumerate}[itemsep=0.2em,topsep=0.2em]
    \item Stack all training-set 3D points into one big point cloud.
    \item Fit \texttt{KMeans(n\_clusters=K)} on this cloud.
    \item For every gesture (train and test), replace each 3D point
          by the index of the nearest centroid.
\end{enumerate}

This rigorous per-fold fitting is \emph{stricter} than the
project guidelines require (which allow one global $k$-means fit). We
keep it because the cost is negligible and it removes a small data
leakage from the report's threat model. The cluster $K$ is chosen by a
\emph{validation curve} on Edit-Distance UI accuracy
(Sec.~\ref{sec:vc}).

% =========================================================================
\section{Dynamic Time Warping (DTW)}

DTW measures the elastic distance between two 3D paths
$\mathbf{s}_1,\mathbf{s}_2$ by aligning them through a warping path
\cite{sakoe1978dtw}. Recurrence:
\[
g[i,j] = \min \begin{cases}
    g[i{-}1,\;j]\;+\;d_{ij} \\
    g[i{-}1,\;j{-}1]\;+\;2 d_{ij} \\
    g[i,\;j{-}1]\;+\;d_{ij}
\end{cases},
\qquad
d_{ij}=\|\mathbf{s}_1[i]-\mathbf{s}_2[j]\|_2.
\]
The diagonal step is weighted by~2 (so that 2 unit-cost diagonal moves
match 2 unit-cost orthogonal moves, removing the bias toward elongated
paths). The final distance is normalised by $I+J$ (so that long
gestures are not penalised systematically), following
\cite{mueller2007info}.

\begin{whybox}{numba JIT compilation}
The DTW recurrence and the Edit recurrence are both annotated
\texttt{@njit(cache=True)}. Numba compiles them to a tight machine-code
loop on first call; the cost is a one-time JIT warm-up done in the
main script. Without numba, 1000 DTW computations per fold $\times$ 10
folds becomes a multi-minute bottleneck on pure Python.
\end{whybox}

% =========================================================================
\section{Edit Distance}

Levenshtein distance \cite{levenshtein1966edit}, computed in $O(L_1
L_2)$ time with the standard Wagner--Fischer DP recurrence
\cite{wagner1974string}:
\[
\text{mat}[i,j]=\min\begin{cases}
    \text{mat}[i{-}1,\;j]+1\\
    \text{mat}[i,\;j{-}1]+1\\
    \text{mat}[i{-}1,\;j{-}1]+ \mathbb{1}\!\bigl[s_1[i{-}1]\ne s_2[j{-}1]\bigr]
\end{cases}.
\]
Operates on the $k$-means symbol streams produced upstream.

% =========================================================================
\section{$k$-Nearest-Neighbours classifier}\label{sec:knn}

DTW and Edit Distance only produce a distance. To turn them into a
classifier we use 1-NN over the training set:
\[
\hat{y}(\mathbf{x}) \;=\; y_{\,j^\star}, \qquad
j^\star = \arg\min_{j\in\text{train}} \mathrm{dist}(\mathbf{x},\mathbf{x}_j).
\]
For the user-dependent setting, the nearest neighbour is computed only
within gestures of the same user.

\subsection{Why we force $k=1$ in the entire downstream pipeline}

Our code produces a validation curve scanning
$k\in\{1,3,5,7,9\}$ for DTW (\texttt{d\{1,4\}\_vc\_knn.png}). The
curve is plotted and saved, but \emph{not} used to pick $k$. The
pipeline forces $k=1$ regardless.

\begin{whybox}{1-NN is the canonical gesture-recognition baseline}
\begin{itemize}[itemsep=0.1em,topsep=0.2em]
\item \cite{wobbrock2007dollar1} originally report the \$1 Recognizer as
      a single-template-match algorithm (implicit 1-NN over preprocessed
      templates).
\item \cite{mezari2018commodity} use 1-NN on top of a Kratz \& Rohs-style
      preprocessing for their commodity-device recognizer.
\item \cite{liu2009uwave} (uWave) use a single nearest template under
      DTW, i.e.\ 1-NN.
\item \cite{mitra2007survey} discuss 1-NN as the standard non-parametric
      baseline for gesture recognition (Section 5).
\end{itemize}
Going to $k>1$ would risk averaging out across same-user-different-style
exemplars and would no longer be the comparable-to-literature 1-NN
baseline. The validation curve remains in the report as transparency.
\end{whybox}

% =========================================================================
\section{\texorpdfstring{\$1}{\$1} Recognizer adapted to 3D
                       (Kratz \& Rohs, 2010)}\label{sec:dollar-raw}

The original \$1 Recognizer \cite{wobbrock2007dollar1} is 2D-only.
We follow the canonical 3D extension of
\cite{kratz2010dollar3}, \emph{A \$3 Gesture Recognizer}.

\begin{whybox}{le \$1 est évalué sur des données BRUTES uniquement}
Le pipeline du \$1 contient sa propre normalisation interne
(resample, translation au centroïde, rotation par axe indicatif,
scaling cube uniforme \cite{wobbrock2007dollar1,kratz2010dollar3}).
Lui fournir des coordonnées \emph{externe-standardisées} per-axe
altère la géométrie utilisée par la rotation cross-product et est
méthodologiquement incohérent. Notre ablation calcule quand même les
3 conditions pour transparence, mais l'évaluation principale du \$1
est forcée à \emph{(a) Aucun pré-traitement}. Sur Domain~4,
l'ablation montre que la standardisation externe améliore
\emph{numériquement} le \$1 ; ce gain est un artefact méthodologique
mentionné dans les limitations du rapport, pas dans les comparaisons
canoniques.
\end{whybox}

\subsection{Step 1 --- resample to $N$ equidistant points}
Cumulative arc-length resampling to $N=64$ equally-spaced 3D points.

\subsection{Step 2 --- translate centroid to origin}
$\mathbf{p}\leftarrow \mathbf{p}-\bar{\mathbf{p}}$.

\subsection{Step 3 --- rotate so $\mathbf{p}_0$ aligns with the centroid}
This is the most subtle 3D extension. The 2D \$1 algorithm uses
$\arctan 2$ in the $xy$-plane. In 3D, $\arctan 2$ is undefined, so we
use the axis-angle formulation:
\[
\theta = \arccos\!\left(\frac{\mathbf{p}_0\cdot\bar{\mathbf{p}}}
                              {\|\mathbf{p}_0\|\,\|\bar{\mathbf{p}}\|}\right),
\qquad
\mathbf{v}_{\text{axis}} =
\frac{\mathbf{p}_0\times\bar{\mathbf{p}}}
     {\|\mathbf{p}_0\times\bar{\mathbf{p}}\|}.
\]
The rotation is then applied via Rodrigues' formula:
\[
\mathbf{v}_{\text{rot}}
 = \mathbf{v}\cos\theta
 + (\mathbf{k}\times\mathbf{v})\sin\theta
 + \mathbf{k}(\mathbf{k}\!\cdot\!\mathbf{v})(1-\cos\theta).
\]
Degenerate cases ($\mathbf{p}_0$ collinear with $\bar{\mathbf{p}}$,
zero norms, $\theta < 10^{-9}$) fall back to identity.

\subsection{Step 4 --- uniform scaling inside a normalised cube of
                       side $l=1$}
$\mathbf{p}\leftarrow \mathbf{p}\,l/\max\bigl(\text{bounding-box edges}\bigr)$.
We do \emph{not} scale axis-by-axis, because for quasi-planar Domain~1
gestures the third extent is near zero and a per-axis division by
that small number is numerically unsafe \cite{kratz2010dollar3}.

\subsection{Confidence score (3D version)}
\[
S \;=\; 1 - \frac{d}{0.5 \,\sqrt{3}\, l},
\qquad
d = \frac{1}{N}\sum_{i=1}^{N}\|\mathbf{c}_i - \mathbf{t}_i\|_2.
\]
The factor $\sqrt{3}$ replaces $\sqrt{2}$ in the 2D paper (the maximum
possible mean point-to-point distance in a cube of side $l$ is half
its diagonal $\sqrt{3}\,l$).

\subsection{Cached template preprocessing}
\cite{wobbrock2007dollar1}: \emph{``For gestures serving as templates,
Steps 1--3 should be carried out once on the raw input points.''} We
preprocess every gesture in the dataset exactly once at the start of
the cross-validation routine, then re-use the cache for every fold.

\subsection{What we did not implement}
Golden-Section-Search over the three Euler angles
\cite{kratz2010dollar3}: would refine the rotation step but multiplies
runtime by $\sim$11 iterations per pair. Marginal accuracy gain on
this dataset.

% =========================================================================
\section{Hand-crafted features for DT, RF, SVM, LR}

\texttt{extract\_features} produces a 50-dim vector (53 with the 3 EVR
appended) summarising each gesture:
\begin{itemize}[itemsep=0.1em,topsep=0.2em]
\item Per-axis: mean, std, min, max, range, skew, kurt (21).
\item Kinematics: speed mean/std/min/max, accel mean/std, arc-length (7).
\item Bounding-box edges and diagonal (4).
\item 3D curvature angles: mean/std/total (3).
\item Per-segment (1/3, 2/3, 3/3) mean speed and displacement (6).
\item Per-axis total path length (3).
\item PCA EVR (3, only when ``Std + PCA denoise'' preprocessing).
\end{itemize}

\begin{whybox}{handcrafted features instead of raw signals}
RF/DT/SVM expect a fixed-length vector. Padding the raw $T\times 3$
sequences to a fixed length blows up the feature space and adds the
implicit assumption that timestep $t$ has the same meaning across
gestures (which it does not). Handcrafted features collapse the
sequence into descriptors that are themselves time-invariant
(Mitra \& Acharya, 2007, \cite{mitra2007survey}, §5).
\end{whybox}

% =========================================================================
\section{Decision Tree}

\cite{quinlan1986id3, breiman1984cart}. GridSearchCV over
\verb|max_depth|, \verb|min_samples_split|, \verb|criterion|
(\verb|gini|/\verb|entropy|), inner \textbf{5-fold} CV
(cf.\ Sec.~\ref{sec:innercv}).

\begin{whybox}{a single tree as a pedagogical control}
We do not expect DT to beat RF. Including it is a controlled experiment
to \emph{show} on our own data the bagging gain
\cite{breiman2001randomforests, hastie2009elements}. If RF beats DT
significantly under Wilcoxon, we have an empirical justification for
the cost of using RF in the deployed system.
\end{whybox}

% =========================================================================
\section{Random Forest}

\cite{breiman2001randomforests}. Bagging of decision trees with
random feature subsampling at each split. We set
\verb|max_features="sqrt"| (the Breiman default for classification).
GridSearchCV over \verb|n_estimators|, \verb|max_depth|,
\verb|min_samples_split| with \textbf{inner 5-fold} CV
(Sec.~\ref{sec:innercv}).

\subsection{Per-fold three-stage feature selection}\label{sec:perfold-fs}

The feature set is large (50--53 dimensions) and contains redundant
descriptors (e.g.\ \texttt{x\_range} $= x_{\max} - x_{\min}$,
\texttt{bbox\_diag} derived from \texttt{bbox\_x/y/z}).  The previous
version computed permutation importance on the \emph{same} training data
used to fit the quick RF --- a methodological flaw that caused the RF to
overestimate the importance of memorised-noise features while
underestimating that of correlated but genuinely useful ones.

\textbf{Current procedure} (applied identically to DT, RF, SVM, LR in
both UI and UD --- function \texttt{\_select\_features\_per\_fold}):

\begin{enumerate}[itemsep=0.15em,topsep=0.2em]
    \item \textbf{Variance filter.}  Remove near-constant features
          ($\mathrm{Var} < 10^{-10}$).  Near-constant features carry
          no information and destabilise Pearson correlation estimates
          \cite{guyon2003fs}.

    \item \textbf{Correlation-based redundancy removal}
          \cite{yu2004relevance,hall1999fs}.  For each pair of features
          with $|r_\text{Pearson}| > 0.95$, remove the lower-variance
          member.  \emph{Why this matters:} if feature $A$ and feature
          $B$ share the same information, permuting $A$ leaves $B$
          intact and the model can recover from the perturbation; $A$
          thus appears \emph{unimportant} even though it is genuinely
          informative \cite{strobl2007bias}.  Without deduplication the
          cumulative-importance cut discards entire correlated clusters,
          destroying classifier performance.

    \item \textbf{Permutation importance on a held-out 20\,\% split.}
          The remaining training samples are split 80/20.  A quick RF
          ($n_\text{estimators}=100$) is fitted on the 80\,\% part; then
          \verb|permutation_importance| (10 repeats) is evaluated on the
          20\,\% held-out part \cite{breiman2001randomforests,
          strobl2007bias}.  Using held-out data gives unbiased importance
          estimates: a model that memorises training noise will not benefit
          from having that noise preserved in the held-out set.  Negative
          importances (feature hurts accuracy) are clipped to~0.  The
          smallest set whose \textbf{cumulative importance} reaches
          90\,\% is retained \cite{guyon2003fs}.
\end{enumerate}

\begin{whybox}{pourquoi trois étapes}
\begin{itemize}[itemsep=0.1em,topsep=0.2em]
\item Gini importance (MDI) is biased toward continuous and
      high-cardinality features \cite{strobl2007bias}; permutation
      importance on held-out data is not.
\item Correlation deduplication is necessary because permutation
      importance is not reliable in the presence of correlated features
      \cite{strobl2007bias,yu2004relevance}: without it, the importance
      of each member of a correlated pair is artificially diluted.
\item Computing importance on the training set itself is biased
      (memorised features look important); a held-out split removes that
      bias \cite{breiman2001randomforests,ambroise2002selection}.
\item All steps use only the outer-fold training set
      \cite{ambroise2002selection}: no test label leaks into the
      feature-selection process.
\end{itemize}
\end{whybox}

\textbf{What this replaces in the report.} Whenever a previous draft
mentioned ``Gini importance threshold $\geq 0.01$ applied in UD only'',
read instead: ``per-fold three-stage feature selection (variance filter
$\to$ correlation deduplication $\to$ permutation importance on
held-out split, cumulative threshold 90\,\%), applied uniformly to
DT/RF/SVM/LR in both UI and UD.''

% =========================================================================
\section{SVM with RBF kernel}

\cite{cortes1995svm, scholkopf2002learning}. The RBF kernel
\[
K(\mathbf{x},\mathbf{y}) = \exp\bigl(-\gamma\,\|\mathbf{x}-\mathbf{y}\|^2\bigr)
\]
maps the 50-dim feature space implicitly into a Hilbert space where
classes become more separable.  Hyperparameters $C$ and $\gamma$
selected by GridSearchCV ($4{\times}4$ grid, inner \textbf{5-fold} CV,
Sec.~\ref{sec:innercv}).

\textbf{Feature scaling.}  SVM with RBF kernel is sensitive to
feature magnitudes: a feature with large range dominates the Euclidean
distance in the kernel, masking the contribution of smaller features.
We apply \texttt{StandardScaler} inside a \texttt{sklearn.Pipeline},
so that the scaler is fitted on each inner-CV training split and never
sees held-out data \cite{hsuchang2010guide}.  This is the canonical
pre-processing recommendation for RBF-SVM:
\begin{quote}
\emph{``We recommend normalising each attribute to $[{-}1,+1]$ or
$[0,1]$''} --- \cite{hsuchang2010guide}, Section 2.
\end{quote}

\begin{whybox}{SVM is a core gesture-recognition reference}
\begin{itemize}[itemsep=0.1em,topsep=0.2em]
\item \cite{mitra2007survey} survey gesture recognition and devote
      §5.4 to SVM-based classifiers.
\item \cite{wuhuang1999vision} review vision-based gesture
      recognition and list SVM as a standard non-temporal classifier.
\item \cite{bhattacharya2014kinect} apply SVM successfully to
      Kinect-derived gesture features.
\end{itemize}
Adding SVM gives our comparison a non-tree-based reference of strong
literature pedigree, complementing the tree-based pair (DT, RF) and the
template-based \$1.
\end{whybox}


\section{Logistic Regression (multinomial, L2)}

\cite{cox1958regression, hosmer2013logistic, hastie2009elements}.
Multinomial L2-penalised LR via the \verb|lbfgs| solver
(\verb|sklearn.linear_model.LogisticRegression|, \verb|max_iter=5000|).
GridSearchCV over the inverse regularisation strength
$C\in\{0.01, 0.1, 1, 10, 100\}$ with \textbf{5-fold} inner CV.

\textbf{Feature scaling.}  The \verb|lbfgs| optimiser requires
inputs on a comparable scale for reliable convergence (the gradient
magnitude is dominated by the largest-scale feature otherwise).
We use \texttt{StandardScaler} inside a \texttt{Pipeline}, fitted on
each inner-CV training split only.  \verb|max_iter=5000| guarantees
convergence for all grid cells; the previous value of 2000 triggered
a \texttt{ConvergenceWarning} on unscaled inputs.

\begin{whybox}{a \emph{linear} classifier as a lower bound}
LR is the simplest non-trivial classifier on a feature vector. Comparing
LR to RBF-SVM directly quantifies how much \emph{non-linear} structure
exists in the feature space: if SVM~$\gg$~LR, the kernel matters; if
SVM~$\approx$~LR, the gain of kernel methods on our features is marginal.
LR also acts as a sanity check: a healthy ML pipeline should produce
a non-trivial LR result (well above chance, 10\,\%).
\end{whybox}

\textbf{In gesture-recognition}, LR is mentioned in
\cite{wuhuang1999vision} (Sec.~3) as a linear baseline. It is rarely
the \emph{best} classifier reported but consistently shows up as a
reference.

% =========================================================================
\section{Cross-validation}

\subsection{User-independent (UI) --- leave-one-user-out}
For each user $u\in\{0,\dots,9\}$: training on the other 9 users
(900 sequences), testing on user $u$ (100 sequences). 10 folds.

\subsection{User-dependent (UD) --- leave-one-sample-out}
For each user, hold out one repetition per gesture as the test fold,
keep the remaining 9. 10 folds; the test fold of fold $k$ contains
exactly 100 sequences (10 gestures $\times$ 10 users $\times$ 1
held-out repetition).

\subsection{Nested CV for tunable classifiers}\label{sec:innercv}
For DT, RF, SVM, LR we use GridSearchCV with an inner \textbf{5-fold}
CV \emph{inside} the outer train fold; no test data ever leaks into the
hyperparameter search.

\begin{whybox}{pourquoi 5-fold (et pas 3-fold) ?}
\cite{varma2006bias} et \cite{cawley2010overfitting} montrent que
3-fold inner CV introduit un biais significatif dans l'estimation
d'erreur pour les datasets de taille modeste. La recommandation
standard est 5--10 folds. Avec $\sim 900$ samples par train fold
externe et 10 classes, 5-fold laisse $\sim 180$ samples par inner
test fold (suffisant pour estimer l'accuracy à $\pm 2\%$).
\end{whybox}

\begin{whybox}{the asymmetry of hyperparameter tuning is on purpose}
DTW, Edit, \$1 have one or two scalar hyperparameters with a clear
plateau, so a single empirical validation-curve scan suffices.
DT/RF/SVM/LR have a combinatorial grid that genuinely needs per-fold
tuning to avoid overfitting in the small-data regime. The comparison
in the final tables is therefore between \emph{best-configured}
versions of each method, not between methods sharing an identical
tuning protocol --- this is acknowledged in the report.
\end{whybox}

% =========================================================================
\section{Hyperparameter validation curves}\label{sec:vc}

We produce two curves per domain:

\subsection{$K_{\text{clusters}}$ scan (Edit Distance UI accuracy)}
Range $\{5,10,15,20,25,30,40\}$. The $K$ that maximises the mean UI
accuracy is the one used in every downstream Edit-Distance evaluation
(and in the ablation study). Saved as
\texttt{d\{1,4\}\_vc\_kclusters.png}.

\subsection{$K_{\text{NN}}$ scan (DTW UI accuracy)}
Range $\{1,3,5,7,9\}$. The curve is plotted and saved
(\texttt{d\{1,4\}\_vc\_knn.png}) but its outcome is \emph{ignored}:
the pipeline forces $k=1$ everywhere. See Sec.~\ref{sec:knn} for the
justification.

% =========================================================================
\section{Ablation study (7 methods $\times$ 3 preprocessings, run in UI \emph{and} UD)}

Three preprocessing conditions:

\begin{enumerate}[label=(\alph*),itemsep=0.2em,topsep=0.2em]
\item No preprocessing (raw 3D coordinates).
\item Per-gesture standardisation.
\item Per-gesture standardisation + PCA denoising 3D$\to$2D$\to$3D
      (the full pipeline). For tree-/kernel-/LR-based classifiers, the
      3 EVR scalars are appended as additional features.
\end{enumerate}

The function \texttt{run\_ablation\_study(setting=...)} now runs in both
\textbf{user-independent} (UI) and \textbf{user-dependent} (UD)
settings. Each setting selects its own best preprocessing per method.

\begin{whybox}{pourquoi l'ablation aussi en UD}
La précédente version supposait implicitement que ``le meilleur
pré-traitement pour UI est aussi le meilleur pour UD''. Cette
hypothèse n'est pas justifiée : un classifieur peut bénéficier de la
PCA-denoising en UI (qui supprime du bruit inter-utilisateur) sans en
bénéficier en UD (où la personnalisation rend ce bruit non pertinent).
On vérifie maintenant empiriquement les divergences.
\end{whybox}

\textbf{Politique \$1 (override post-ablation)} : après l'ablation,
\$1 est forcé à condition (a) raw quel que soit l'optimum trouvé
(cf.\ Sec.~\ref{sec:dollar-raw}). L'override est appliqué identiquement
pour UI et UD.

\textbf{Outputs} : 4 CSVs au format
\texttt{Outputs/tables/ablation/ablation\_domain\{1,4\}\_\{ui,ud\}.csv}
(21 lignes chacun : 7 méthodes $\times$ 3 conditions). Plus un fichier
\texttt{preproc\_ui\_vs\_ud\_domain\{1,4\}.csv} listant, par méthode, le
best-preproc UI vs.\ UD avec un flag \verb|Divergent|.

% =========================================================================
\section{Statistical tests}

\subsection{Per-(gesture, user) accuracy vector ($n=100$)}
For every method, we keep a 100-dim accuracy vector where entry
$(g, u)$ is the mean correctness over the test-fold repetitions of
gesture~$g$ by user~$u$.

\subsection{Wilcoxon signed-rank test}
For every pair of methods (M_A, M_B), we run a paired Wilcoxon
signed-rank test on their 100-vectors. With \textbf{7 methods} we have
$\binom{7}{2}=21$ pairs per domain.

\subsection{Degenerate case}
If all paired differences are zero (two methods identical on every
$(g,u)$), \verb|scipy.stats.wilcoxon| raises. The wrapper
\verb|_safe_wilcoxon| returns $p=1$ with a warning.

\subsection{Multiple-comparison corrections}
\textbf{Bonferroni} ($\alpha/n_{\text{comp}}$, conservative) and
\textbf{Benjamini-Hochberg FDR} (more powerful, controls false
discovery rate at $5\%$). Both are reported side-by-side. The
permutation test and the Bayesian sign test that were in earlier
drafts have been removed: they are redundant with Wilcoxon~+~correction.

Saved as \texttt{Outputs/tables/statistical\_tests/p\_values\_domain\{1,4\}\_user\_independent.csv}
(symmetric $6\times 6$ matrix) and a heatmap PNG.

% =========================================================================
\section{Overfitting diagnostic}\label{sec:overfitting}

Le prof a fait remarquer qu'une accuracy proche de 100\,\% (comme le
\$1 UD D1 à 0.998) est suspecte et peut cacher du sur-apprentissage.
On ajoute donc deux diagnostics complémentaires.

\subsection{Gap train vs.\ test par fold}

Pour chaque \emph{classifieur paramétrique} (DT, RF, SVM, LR), les
fonctions \texttt{crossval\_\{ui,ud\}\_*} retournent désormais un
5\textsuperscript{e} élément : la liste des accuracies train, fold par
fold. Le rapport produit \texttt{overfitting\_gap.csv} avec :
\[
\text{Gap} \;=\; \overline{\text{train\_acc}} \;-\; \overline{\text{test\_acc}}.
\]
Un Gap $> 0.05$ est flaggé comme « high variance ».

\textbf{Les méthodes 1-NN (DTW, Edit, \$1) sont exclues} de cette
analyse : leur train accuracy est 1.0 par construction (chaque sample
d'entraînement est son propre plus proche voisin), donc le gap n'a pas
de sens pour elles. La quasi-perfection observée pour \$1 en UD D1
\textbf{n'est pas} de l'overfitting au sens machine-learning ; c'est
simplement que la personnalisation rend le 1-NN trivialement résoluble
sur ce domaine planaire.

\subsection{Learning curves \texttt{sklearn.learning\_curve}}

Pour chaque classifieur paramétrique et chaque domaine, on trace la
courbe d'apprentissage : accuracy train et accuracy validation en
fonction de la taille du training set
$\{0.2, 0.4, 0.6, 0.8, 1.0\} \times N_\text{train}$, avec 5-fold CV.
Référence : Hastie et al.\ \cite{hastie2009elements}, §7.10.

\begin{whybox}{lecture biais-variance des courbes d'apprentissage}
\begin{itemize}[itemsep=0.1em,topsep=0.2em]
\item \textbf{Écart train$-$val large et persistant} = high variance =
      overfitting. Remède : régulariser, ajouter des données,
      simplifier le modèle.
\item \textbf{Train et val tous deux bas} = high bias = underfitting.
      Remède : modèle plus capable, plus de features.
\item \textbf{Train et val proches et élevés, plateau atteint} =
      modèle bien calibré pour ce dataset.
\end{itemize}
\cite{domingos2012few} : \emph{``the most important factor by far is
the features used.''} Si on observe high bias, c'est la feature engineering
qu'il faut revoir.
\end{whybox}

\textbf{Outputs} : \texttt{Outputs/figures/learning\_curves/lc\_\{dt,rf,svm,lr\}\_d\{1,4\}\_ui.png}
(8 figures) + \texttt{Outputs/tables/overfitting/overfitting\_gap.csv}.


\section{Confusion matrix of the best model}

After the statistical tests, we pick the method with the highest mean
UI accuracy for each domain, refit it across the 10 UI folds, and plot
the aggregated confusion matrix
(\texttt{Outputs/figures/confusion\_matrices/}).

\begin{whybox}{what to look at in the confusion matrix}
\textbf{Off-diagonal hot-spots} = pairs of gestures the model
confuses (e.g.\ ``2'' versus ``5'' in Domain~1 because they share
similar arcs; cone vs.\ pyramid in Domain~4 because the only
distinguishing feature is curvature near the base). Discussing
these explicitly in the oral defence is more interesting than
quoting global accuracy numbers.
\end{whybox}

% =========================================================================
\section{Output organisation}\label{sec:outputs}

\begin{lstlisting}[caption={Outputs/ tree as produced by CODE\_IA.py}]
Outputs/
|- figures/
|   |- exploratory/             sequence_lengths, gesture_samples
|   |- pca_denoising/           per-domain EVR distributions
|   |- validation_curves/       K_CLUSTERS and KNN_K scans
|   |- feature_importance/      RF permutation-importance bar plots
|   |- confusion_matrices/      best-model CMs (one per domain)
|   |- statistical_tests/       Wilcoxon p-value heatmaps
|   `- learning_curves/         lc_{dt,rf,svm,lr}_d{1,4}_ui.png
|- tables/
|   |- ablation/                ablation_domain{1,4}_{ui,ud}.csv
|   |                           + preproc_ui_vs_ud_domain{1,4}.csv
|   |- fold_results/            per-method per-fold accuracies
|   |- statistical_tests/       p_values_domain{1,4}_*.csv
|   |- summary/                 summary_results.csv
|   `- overfitting/             overfitting_gap.csv
`- Documentation/
    |- pipeline_explained.{tex,pdf}    THIS DOCUMENT
    `- results_synthesis.{tex,pdf}     companion results memo
\end{lstlisting}

% =========================================================================
\section{Limitations and choices not implemented}

\begin{itemize}[itemsep=0.2em,topsep=0.2em]
\item \textbf{PCA / $k$-means data leakage} --- explicitly allowed by
      the project guidelines ``for simplicity''. We chose to fit
      $k$-means \emph{per fold} anyway, but PCA is fit per-gesture
      (which is leak-free by construction).
\item \textbf{Feature selection RF-based filter} --- the per-fold
      permutation importance is computed using a quick RF and applied
      to DT, RF, SVM, LR alike. A method-agnostic filter (e.g.\ mutual
      information) or method-specific filters would be more rigorous
      but multiplies cost. Mentioned as future work.
\item \textbf{Auto feature-extraction libraries} (tsfresh
      \cite{christ2018tsfresh}, tsfel \cite{barandas2020tsfel}) were
      considered but not used: our hand-crafted feature set was
      already large (50--53 features), and the per-fold permutation
      importance plus 90\,\% cumulative threshold gives us empirical
      pruning. Discussed in the report.
\item \textbf{Golden Section Search} in the \$1 pipeline
      \cite{kratz2010dollar3}: not implemented because of its
      $\sim 11\times$ runtime cost for marginal accuracy. Mentioned
      in the report's ``future work''.
\item \textbf{LSTM} was removed earlier on the professor's advice:
      $\sim 1000$ training sequences are too few to justify a
      recurrent network without massive augmentation.
\item \textbf{Per-fold $k$-NN $K$ tuning} replaced by enforced
      $k=1$ (see Sec.~\ref{sec:knn}).
\item \textbf{\$1 on standardised data} --- the ablation reports the
      raw $\to$ std $\to$ PCA conditions for transparency, but the main
      evaluation forces (a) raw for \$1 per Wobbrock 2007 / Kratz \&
      Rohs 2010. The reader interested in the (b)/(c) numbers should
      look at the ablation CSV; they are not part of the canonical
      comparison.
\end{itemize}

% =========================================================================
\section*{References}
\addcontentsline{toc}{section}{References}
\renewcommand{\refname}{}

\begin{thebibliography}{99}

\bibitem{ambroise2002selection}
Ambroise, C., \& McLachlan, G.\ (2002). Selection bias in gene
extraction on the basis of microarray gene-expression data.
\textit{PNAS}, 99(10), 6562--6566.

\bibitem{barandas2020tsfel}
Barandas, M., et al.\ (2020). TSFEL: Time series feature extraction
library. \textit{SoftwareX}, 11, 100456.

\bibitem{breiman1984cart}
Breiman, L., Friedman, J., Olshen, R., \& Stone, C.\ (1984).
\textit{Classification and Regression Trees}. Wadsworth.

\bibitem{breiman2001randomforests}
Breiman, L.\ (2001). Random forests.
\textit{Machine Learning}, 45(1), 5--32.

\bibitem{bhattacharya2014kinect}
Bhattacharya, S., Czejdo, B., \& Perez, N.\ (2014).
Gesture classification with machine learning using Kinect sensor data.
\textit{ICETET'14}.

\bibitem{cawley2010overfitting}
Cawley, G., \& Talbot, N.\ (2010). On over-fitting in model selection
and subsequent selection bias in performance evaluation.
\textit{JMLR}, 11, 2079--2107.

\bibitem{christ2018tsfresh}
Christ, M., Braun, N., Neuffer, J., \& Kempa-Liehr, A.~W.\ (2018).
Time series feature extraction on basis of scalable hypothesis tests
(tsfresh). \textit{Neurocomputing}, 307, 72--77.

\bibitem{cortes1995svm}
Cortes, C., \& Vapnik, V.\ (1995). Support-vector networks.
\textit{Machine Learning}, 20, 273--297.

\bibitem{coverhart1967nn}
Cover, T., \& Hart, P.\ (1967). Nearest neighbor pattern
classification. \textit{IEEE Trans.\ Info.\ Theory}, 13(1), 21--27.

\bibitem{cox1958regression}
Cox, D.~R.\ (1958). The regression analysis of binary sequences.
\textit{J.\ Royal Statistical Society B}, 20(2), 215--242.

\bibitem{domingos2012few}
Domingos, P.\ (2012). A few useful things to know about machine
learning. \textit{Communications of the ACM}, 55(10), 78--87.

\bibitem{guyon2003fs}
Guyon, I., \& Elisseeff, A.\ (2003). An introduction to variable and
feature selection. \textit{JMLR}, 3, 1157--1182.

\bibitem{hosmer2013logistic}
Hosmer, D.~W., Lemeshow, S., \& Sturdivant, R.~X.\ (2013).
\textit{Applied Logistic Regression} (3rd ed.). Wiley.

\bibitem{hall1999fs}
Hall, M.~A.\ (1999). \textit{Correlation-based Feature Selection for
Machine Learning}. PhD thesis, University of Waikato.

\bibitem{hsuchang2010guide}
Hsu, C.-W., Chang, C.-C., \& Lin, C.-J.\ (2010). A practical guide to
support vector classification. Technical report, National Taiwan
University. \texttt{https://www.csie.ntu.edu.tw/{\textasciitilde}cjlin/papers/guide/guide.pdf}

\bibitem{strobl2007bias}
Strobl, C., Boulesteix, A.-L., Zeileis, A., \& Hothorn, T.\ (2007).
Bias in random forest variable importance measures: Illustrations,
sources and a solution. \textit{BMC Bioinformatics}, 8(25).

\bibitem{yu2004relevance}
Yu, L., \& Liu, H.\ (2004). Efficient feature selection via analysis of
relevance and redundancy. \textit{JMLR}, 5, 1205--1224.

\bibitem{varma2006bias}
Varma, S., \& Simon, R.\ (2006). Bias in error estimation when using
cross-validation for model selection.
\textit{BMC Bioinformatics}, 7(91).

\bibitem{hastie2009elements}
Hastie, T., Tibshirani, R., \& Friedman, J.\ (2009).
\textit{The Elements of Statistical Learning} (2nd ed.). Springer.
(Chapters 9, 15.)

\bibitem{huang2019gesture}
Huang, Jaiswal, \& Rai (2019). Gesture-based system for next
generation natural and intuitive interfaces.
\textit{AI EDAM}, 33(1), 54--68.

\bibitem{jolliffe2002pca}
Jolliffe, I.\ (2002). \textit{Principal Component Analysis}
(2nd ed.). Springer.

\bibitem{kratz2010dollar3}
Kratz, S., \& Rohs, M.\ (2010). A \$3 gesture recognizer: simple
gesture recognition for devices equipped with 3D acceleration
sensors. \textit{IUI'10}, 341--344.

\bibitem{levenshtein1966edit}
Levenshtein, V.\ (1966). Binary codes capable of correcting deletions,
insertions, and reversals. \textit{Soviet Physics Doklady},
10, 707--710.

\bibitem{liu2009uwave}
Liu, J., Zhong, L., Wickramasuriya, J., \& Vasudevan, V.\ (2009).
uWave: Accelerometer-based personalized gesture recognition.
\textit{Pervasive and Mobile Computing}, 5(6), 657--675.

\bibitem{mezari2018commodity}
Mezari, A., \& Maglogiannis, I.\ (2018). An easily customized gesture
recognizer for assisted living using commodity mobile devices.
\textit{J.\ Healthcare Engineering}, 2018:3180652.

\bibitem{mitra2007survey}
Mitra, S., \& Acharya, T.\ (2007). Gesture recognition: a survey.
\textit{IEEE Trans.\ SMC-C}, 37(3), 311--324.

\bibitem{mueller2007info}
M\"uller, M.\ (2007). \textit{Information Retrieval for Music and
Motion}. Springer. (Chapter 4: DTW.)

\bibitem{quinlan1986id3}
Quinlan, J.~R.\ (1986). Induction of decision trees.
\textit{Machine Learning}, 1(1), 81--106.

\bibitem{rabiner1993fundamentals}
Rabiner, L., \& Juang, B.-H.\ (1993). \textit{Fundamentals of Speech
Recognition}. Prentice-Hall.

\bibitem{sakoe1978dtw}
Sakoe, H., \& Chiba, S.\ (1978). Dynamic programming algorithm
optimization for spoken word recognition. \textit{IEEE Trans.\
ASSP}, 26(1), 43--49.

\bibitem{scholkopf2002learning}
Sch\"olkopf, B., \& Smola, A.\ (2002). \textit{Learning with Kernels}.
MIT Press.

\bibitem{theodoridis2009pattern}
Theodoridis, S., \& Koutroumbas, K.\ (2009). \textit{Pattern
Recognition} (4th ed.). Academic Press.

\bibitem{wagner1974string}
Wagner, R., \& Fischer, M.\ (1974). The string-to-string correction
problem. \textit{J.\ ACM}, 21(1), 168--173.

\bibitem{wobbrock2007dollar1}
Wobbrock, J.~O., Wilson, A.~D., \& Li, Y.\ (2007). Gestures without
libraries, toolkits or training: A \$1 recognizer for user interface
prototypes. \textit{UIST'07}, 159--168.

\bibitem{wuhuang1999vision}
Wu, Y., \& Huang, T.~S.\ (1999). Vision-based gesture recognition: a
review. \textit{Gesture Workshop}, LNAI 1739, 103--115.

\end{thebibliography}

\end{document}
